% !TeX root = ../main.tex
\section{Sygnał Wi-Fi IEEE 802.11}
\label{sec:wifi_phy}

% ------------------------------------------------------------------------
% Rodzina standardów IEEE 802.11 i ewolucja warstwy fizycznej (802.11a -> 802.11ax)
% Modulacja OFDM i prefiks cykliczny (zasada działania, okres FFT, rola CP)
% Struktura ramki fizycznej: Non-HT (802.11a) oraz HE-SU (802.11ax)
% Parametry radarowe sygnałów IEEE 802.11a i IEEE 802.11ax:
%   - tabela porównawcza parametrów (pasmo, Nfft, d_f, Tsym, Tcp, dR, v_unamb)
%   - drastyczny spadek prędkości jednoznacznej w 802.11ax (ok. 16.7 m/s vs 68.2 m/s)
% Siatka podnośnych OFDM i problem luki nośnych stałoprądowych (DC):
%   - 1 wyzerowana nośna w 802.11a vs 3 nośne w 802.11ax i ich wpływ na wyciek widmowy
% Analiza transmisji pakietowej (przerwy SIFS/DIFS, limit PSDU) i potencjał Multi-Frame CPI
% ------------------------------------------------------------------------

\subsection{Rodzina standardów IEEE 802.11 i ewolucja warstwy fizycznej}
\label{subsec:wifi_standards}

\subsection{Modulacja OFDM i prefiks cykliczny}
\label{subsec:wifi_ofdm_cp}

\subsection{Struktura ramki fizycznej: Non-HT (802.11a) oraz HE-SU (802.11ax)}
\label{subsec:wifi_frame_structure}

\subsection{Porównanie parametrów radarowych sygnałów IEEE 802.11a i IEEE 802.11ax}
\label{subsec:wifi_radar_parameters}

\subsection{Siatka podnośnych OFDM i problem luki nośnych stałoprądowych (DC)}
\label{subsec:wifi_subcarriers_dc}
